% from ega4-14-fr.tex
\oldpage[IV-3]{201}
on ait $V\cap Y_i\subset\psi(U)\cap Y_i$ et comme la réunion des $V\cap Y_i$ pour
ces indices est $V$, on a $V\subset\psi(U)$, donc $\psi(U)$ est un voisinage de
$\psi(x)$ dans $Y$. L'assertion (iii) est triviale.
\end{proof}

\begin{remarks}[14.1.3]
\label{IV.14.1.3-fr}
\begin{enumerate}[label=(\roman*)]
  \item L'ensemble $Z$ des points $x\in X$ où un morphisme est ouvert
  \emph{n'est pas nécessairement ouvert}. Par exemple, soient $K$ un corps, $A$
  l'anneau de polynômes $K[S,T]$, $V$ le plan affine $\operatorname{Spec}(A)$,
  $X$ le sous-préschéma fermé de $V$ «~réunion de la droite $X_1$ définie par
  $T=0$ et de la droite $X_2$ définie par $S=0$~», c'est-à-dire
  $\operatorname{Spec}(A/\mathfrak a)$, où $\mathfrak a=AST$~; prenons
  $Y=X_2=\operatorname{Spec}(A/AS)$ et pour $f$ la projection correspondant à
  l'injection canonique $K[T]\to A/\mathfrak a$~; alors on a $Z=X_2$, qui n'est
  pas ouvert dans $X$.

  \item Soient $X$, $Y$ deux préschémas noethériens irréductibles de points
  génériques $\xi$, $\eta$ respectivement, $f:X\to Y$ un morphisme
  \emph{dominant localement de type fini}~; alors $f$ \emph{est ouvert au point
  $\xi$}. En effet, on peut évidemment se borner au cas où $X$ et $Y$ sont réduits
  (donc intègres) (1.5.4) et affines~; en vertu du théorème de platitude générique
  (6.9.1), il existe un ouvert non vide $V$ de $Y$ tel que la restriction
  $f^{-1}(V)\to V$ de $f$ soit un morphisme \emph{plat}, et on conclut de (2.4.6)
  que cette restriction est un morphisme ouvert. Toutefois, comme il y a des
  morphismes dominants de type fini $f:X\to Y$ (où $X$ et $Y$ sont
  \emph{irréductibles}) qui ne sont pas ouverts (voir par exemple (\textbf{II},
  8.1)), l'ensemble des points où un tel morphisme est ouvert n'est pas
  nécessairement \emph{fermé} dans $X$.

  \item Nous ignorons si, lorsque $X$ et $Y$ sont localement noethériens et
  $f:X\to Y$ un morphisme localement de type fini, l'ensemble des points de $X$
  où $f$ est ouvert est ou non \emph{localement constructible}.
\end{enumerate}
\end{remarks}

\begin{proposition}[14.1.4]
\label{IV.14.1.4-fr}
Soient $X$, $Y$ deux espaces topologiques, $\psi:X\to Y$ une application
continue. Pour tout $y\in Y$, l'ensemble des $x\in\psi^{-1}(y)$ où $\psi$ est
ouvert est une partie fermée de $\psi^{-1}(y)$.
\end{proposition}

\begin{proof}
En effet, supposons que $\psi$ ne soit pas ouvert en un point
$x\in\psi^{-1}(y)$~; il existe alors un voisinage ouvert $V$ de $x$ dans $X$ tel
que $\psi(V)$ ne soit pas un voisinage de $y$~; il en résulte que pour tout
$x'\in V\cap\psi^{-1}(y)$, $\psi$ n'est pas ouvert au point $x'$.
\end{proof}

\begin{remark}[14.1.5]
\label{IV.14.1.5-fr}
Même si $X$ et $Y$ sont localement noethériens et $f$ un morphisme \emph{fini},
il peut se faire que $f$ soit ouvert en tous les points d'une fibre $f^{-1}(y)$,
sans qu'il existe un \emph{voisinage} de $f^{-1}(y)$ en tous les points duquel $f$
soit ouvert. Soient par exemple $K$ un corps, $A$ l'anneau de polynômes
$K[T_1,T_2,T_3]$, $V=\operatorname{Spec}(A)$ l'espace affine à 3 dimensions sur
$K$, $X=\operatorname{Spec}(A/\mathfrak a)$, où $\mathfrak a=\mathfrak b\mathfrak c$,
avec $\mathfrak b=(T_3)$ et $\mathfrak c=(T_1)+(T_2-T_3)$ dans $A$, de sorte que
$X$ est réunion du plan $X_1=\operatorname{Spec}(A/\mathfrak b)$ («~plan
d'équation $T_3=0$~») et de la droite $X_2=\operatorname{Spec}(A/\mathfrak c)$
(«~droite d'équations $T_1=0$, $T_2=T_3$~») qui sont ses composantes
irréductibles. Prenons $Y=X_1$ et soit $f:X\to Y$ la projection correspondant à
l'injection canonique $K[T_1,T_2]\to A/\mathfrak a$~; si $y$ est le point commun
à $X_1$ et $X_2$, $f^{-1}(y)$ est réduit à $y$ et $f$ est ouvert en ce point mais
n'est ouvert en aucun point de $X_2$ dans un voisinage de $y$ et distinct de $y$.
\end{remark}

\begin{env}[14.1.6]
\label{IV.14.1.6-fr}
Dans ce qui suit, le rôle essentiel sera joué par le critère (1.10.3)
caractérisant les morphismes \emph{localement de présentation finie} $f:X\to Y$
qui sont ouverts en un point $x$ par le «~relèvement des générisations~»~: pour
toute générisation $y'$ de $y=f(x)$, il existe $x'\in X$, générisation de $x$, tel
que $y'=f(x')$.
\end{env}

